\documentclass[12pt]{article}
\usepackage{amssymb}
\usepackage{geometry}
\usepackage{graphicx}
\usepackage[colorlinks=true, allcolors=blue]{hyperref}
\usepackage{cite}
\usepackage{array}
\usepackage{booktabs}
\usepackage{wrapfig}
\usepackage{amsmath}
\usepackage{lipsum}
\usepackage{caption}
\usepackage{ulem}
\usepackage{xcolor}
\usepackage{tikz}
\usetikzlibrary{arrows.meta, positioning, shapes}
\usepackage{sectsty}

\sectionfont{\color{sectionColor}}
\subsectionfont{\color{subsectionColor}}

\definecolor{titleColor}{HTML}{800000}        % Maroon
\definecolor{sectionColor}{HTML}{4682B4}      % SteelBlue
\definecolor{textHighlight}{HTML}{DAA520}     % Goldenrod
\definecolor{backgroundColor}{HTML}{F0F8FF}   % AliceBlue
\definecolor{emphasisColor}{HTML}{228B22}     % ForestGreen
\definecolor{subsectionColor}{HTML}{6A5ACD}   % SlateBlue

\title{\textcolor{titleColor}{\textbf{\Huge From Computers to the Heart--Brain System: Signals and Meaning}}}
\author{B. M. Osman\thanks{Al Nahada University, Cairo \\ \texttt{babikerosman@yahoo.com}}}
\date{\today}

\begin{document}
\maketitle

\section*{Introduction}
At the most fundamental level, all systems of information processing rely on \textit{signals}.
Computers operate on discrete electrical signals \cite{vonNeumann1945, turing1937},
the brain on dynamic electrochemical signals \cite{kandel2000principles, dayan2001theoretical},
and the heart--brain system, as conceptualized in the Heart-Based Resonant Consciousness Theory (HBRCT),
on an integration of electromagnetic and neurochemical signals \cite{mcCraty2003heart, bradley2009electrophysiology}.
The distinction lies not only in the \textbf{nature of the signals}, but also in how these signals are organized to produce \textbf{logic, thought, and meaning}.
Whereas the computer reduces signals to binary computation, and the brain elevates signals to adaptive thought,
the heart--brain system transcends both by embedding coherence, emotional tone, and collective resonance.
\section*{Tabular Comparison}

\begin{tabular}{>{\raggedright\arraybackslash}p{3cm} 
                >{\raggedright\arraybackslash}p{4cm} 
                >{\raggedright\arraybackslash}p{4cm} 
                >{\raggedright\arraybackslash}p{4cm}}
\toprule
\textbf{Aspect} & \textbf{Computer} & \textbf{Brain} & \textbf{Heart--Brain System (HBRCT)} \\
\midrule
Nature of Signals 
& Electrical (binary on/off, 1s and 0s) 
& Electrochemical (neural firing, continuous and adaptive) 
& Electromagnetic and neurochemical (heart rhythms generating coherence) \\
\midrule
Processing Units 
& Transistors and logic gates 
& Neurons and synapses 
& Heart field + brain networks in resonance \\
\midrule
Information Handling 
& Fixed rules, programmed logic 
& Learning and neuroplasticity 
& Meaning, coherence, and resonance (integration of emotion and intuition) \\
\midrule
Source of Purpose 
& External programmer 
& Experience and memory 
& Intrinsic resonance, love, and collective consciousness \\
\midrule
Output 
& Logical results, computations 
& Thoughts, perceptions, adaptive behavior 
& Coherence, intuitive insight, collective field of consciousness \\
\bottomrule
\end{tabular}

\section*{Synthesis}
This progression from \textbf{logic (computer)} to \textbf{thought (brain)} to \textbf{meaning and coherence (heart--brain)} highlights a spectrum of signal organization.
The computer demonstrates the foundation of physical signal dependency,
the brain introduces adaptability and cognition,
and the heart--brain system integrates a higher dimension of coherence,
linking individual consciousness to a collective field.
In this way, HBRCT proposes that the heart is not merely a physiological pump,
but a resonant integrator of signals, encoding meaning, intuition, and love into the dynamics of consciousness \cite{mcCraty2015science, childre2016heartintelligence}.

\section*{Conclusion}
The transition from computer logic to brain thought to heart coherence represents an evolutionary deepening of signal processing: from raw computation, to adaptive cognition, to the emergence of meaning and resonance.
This layered view situates HBRCT as a new framework for unifying biological, computational, and conscious processes.

\nocite{*}
\bibliographystyle{plain}
\bibliography{ref}

\end{document}






\bibliographystyle{unsrturl}


\usepackage[backend=bibtex,style=numeric,sorting=none]{biblatex}
\addbibresource{ref.bib}


\nocite{*}
\printbibliography



\section*{Tabular Comparison}

\begin{tabular}{>{\raggedright\arraybackslash}p{3cm} 
                >{\raggedright\arraybackslash}p{4cm} 
                >{\raggedright\arraybackslash}p{4cm} 
                >{\raggedright\arraybackslash}p{4cm}}
\toprule
\textbf{Aspect} & \textbf{Computer} & \textbf{Brain} & \textbf{Heart--Brain System (HBRCT)} \\
\midrule
Nature of Signals 
& Electrical (binary on/off, 1s and 0s) 
& Electrochemical (neural firing, continuous and adaptive) 
& Electromagnetic and neurochemical (heart rhythms generating coherence) \\
\midrule
Processing Units 
& Transistors and logic gates 
& Neurons and synapses 
& Heart field + brain networks in resonance \\
\midrule
Information Handling 
& Fixed rules, programmed logic 
& Learning and neuroplasticity 
& Meaning, coherence, and resonance (integration of emotion and intuition) \\
\midrule
Source of Purpose 
& External programmer 
& Experience and memory 
& Intrinsic resonance, love, and collective consciousness \\
\midrule
Output 
& Logical results, computations 
& Thoughts, perceptions, adaptive behavior 
& Coherence, intuitive insight, collective field of consciousness \\
\bottomrule
\end{tabular}
